% GLOSARIUM

\chapter*{GLOSARIUM}
\addcontentsline{toc}{chapter}{GLOSARIUM}
\hyphenpenalty=10000
\exhyphenpenalty=10000
\sloppy

\section*{Istilah Teknis}

\subsection*{A}
\begin{description}[leftmargin=3.5cm, style=sameline, font=\normalfont]
    \item[API] Antarmuka pemrograman aplikasi yang memungkinkan pertukaran data antar sistem perangkat lunak.
\end{description}

\subsection*{B}
\begin{description}[leftmargin=3.5cm, style=sameline, font=\normalfont]
    \item[Backend] Bagian sistem yang berjalan di server untuk memproses logika bisnis dan mengelola basis data.
\end{description}

\subsection*{F}
\begin{description}[leftmargin=3.5cm, style=sameline, font=\normalfont]
    \item[Frontend] Bagian sistem yang tampil di sisi pengguna untuk menangani antarmuka dan interaksi.
\end{description}

\subsection*{J}
\begin{description}[leftmargin=3.5cm, style=sameline, font=\normalfont]
    \item[JWT] JSON Web Token, standar untuk mentransmisikan informasi terenkripsi secara aman antar sistem.
\end{description}

\subsection*{R}
\begin{description}[leftmargin=3.5cm, style=sameline, font=\normalfont]
    \item[Realtime] Penyajian atau pemrosesan data secara langsung tanpa jeda waktu signifikan.
    \item[Responsive] Kemampuan tampilan antarmuka untuk menyesuaikan dengan berbagai ukuran layar.
\end{description}

\newpage

\section*{Istilah Non-Teknis}

\subsection*{K}
\begin{description}[leftmargin=3.5cm, style=sameline, font=\normalfont]
    \item[Kampus] Lingkungan atau kompleks bangunan institusi pendidikan tinggi beserta fasilitasnya.
    \item[Kehadiran] Keberadaan atau kedatangan peminjam pada waktu dan tempat yang telah dijadwalkan.
\end{description}

\subsection*{M}
\begin{description}[leftmargin=3.5cm, style=sameline, font=\normalfont]
    \item[Mahasiswa] Peserta didik yang terdaftar dan menempuh pendidikan di perguruan tinggi.
\end{description}

\subsection*{P}
\begin{description}[leftmargin=3.5cm, style=sameline, font=\normalfont]
    \item[Peminjaman] Proses menggunakan fasilitas milik pihak lain untuk jangka waktu tertentu.
    \item[Prasarana] Fasilitas dasar pendukung yang tidak digunakan langsung dalam kegiatan pendidikan.
\end{description}

\subsection*{S}
\begin{description}[leftmargin=3.5cm, style=sameline, font=\normalfont]
    \item[Sarana] Fasilitas yang digunakan secara langsung untuk menunjang kegiatan pendidikan.
    \item[Sarpras] Singkatan dari Sarana dan Prasarana, unit pengelola fasilitas di lingkungan kampus.
    \item[Security] Petugas keamanan yang bertugas menjaga ketertiban dan keamanan lingkungan kampus.
\end{description}

\subsection*{U}
\begin{description}[leftmargin=3.5cm, style=sameline, font=\normalfont]
    \item[UKM] Unit Kegiatan Mahasiswa, organisasi yang menampung minat dan bakat mahasiswa.
\end{description}

\subsection*{V}
\begin{description}[leftmargin=3.5cm, style=sameline, font=\normalfont]
    \item[Verifikasi] Proses pemeriksaan untuk memastikan kebenaran atau kesesuaian suatu data atau informasi.
\end{description}